\subsection{Full-benchmark family-specialized baselines}
\label{sec:leaderboard}

The joint-37 comparison of \cref{sec:mainresults} is a deliberate stress test; it is not
what the benchmark can support under a practical recipe. To calibrate the latter --- and to
give the community a starting leaderboard --- we train each family \emph{separately} on the
biology-CPT encoder init (\cref{sec:archs} footnote: the CPT arm) with the same stabilized
protocol and per-family dev early stopping, for both decision architectures, and evaluate
every supported task's test split (\cref{tab:leaderboard}, single training seed; the
variance caveats of \cref{sec:seedavg} apply and are why the matched comparison above is
seed-averaged).

Three findings. \emph{First, specialization plus CPT recovers near-SOTA numbers}:
TAPE fluorescence reaches Spearman 0.627 with task heads (0.568 with the shared scorer),
against the published ceiling of 0.68; dna\_core reaches accuracy 0.625 (from 0.475 under
joint no-CPT training); and the local-snapshot family reaches 0.860 on nuclear-protein
positioning, 0.855 on signal-peptide detection and 0.785 on promoter detection
(single-task specialized recipes reach 0.88--0.91 on the latter); and Genomic-Benchmarks
proposition tasks reach AUROC 0.897/0.839/0.834 on coding-vs-intergenic,
human-or-worm and non-TATA promoters. Not every family improves: GUE's 28 tasks
average 0.497 accuracy under specialization (best 0.570), and the virus\_covid task
scores 0.115 --- its test split carries 1{,}754 quarantined conflicting-label groups
and an adversarial class balance, which the leaderboard reports rather than hides.
ProteinGym is the starkest headroom: 185 of 217 assays are evaluable within the
512-token budget, and family-specialized training reaches only 0.012 macro Spearman with task
heads ($-0.022$ with the shared scorer; median 0.017, best 0.399) against
MSA-augmented supervised ceilings near 0.55 ---
single-sequence mutation-fitness prediction at this sample and update budget remains
open, and the benchmark quantifies the gap. The low absolute scores of \cref{sec:mainresults} are therefore a
property of the joint small-budget regime and the encoder init, not of the tasks being
unlearnable. \emph{Second, the architecture ordering replicates on the score interface}:
on TAPE stability, per-task heads reach 0.433 while the shared scorer's five-anchor
distribution reaches 0.049 --- the same heads-over-shared-scorer ordering as the seed-averaged
noul result, now on a different primitive and a different training regime, which argues
against the main comparison being an artifact of joint training. \emph{Third, the benchmark
is not saturated}: DeepSTARR (0.23/0.21 vs.\ published PCC 0.68/0.74; our 512-token window
truncates its ${\sim}$2\,kb inputs), the legacy splice/TF blind tests (constant-prediction
collapse in both arms), and several ProteinGym assays remain far below published bests
(\cref{tab:priorbest}), so the benchmark discriminates and leaves measurable headroom.

\paragraph{The cost of unification.}
One checkpoint trained on all 233 gated tasks (\cref{tab:leaderboard}, ``single'' columns)
serves classification and proposition families nearly losslessly (Genomic-Benchmarks noul
AUROC 0.698 vs 0.715 specialized) but sacrifices small regression families (TAPE 0.299 vs
0.530); the heads-over-shared-scorer ordering holds here as well (0.698 vs 0.513), the
fourth regime in which the main comparison replicates.

\paragraph{Granularity gradient.}
Five tasks fine-tuned individually (\cref{fig:gradient}) reach 0.845 promoter, 0.671
fluorescence (99\% of the published ceiling) and 0.892 coding AUROC; per-task ordering
across the four granularities shows within-family positive transfer (family-joint $\ge$
single for npp/fold/coding) and cross-family interference (single $>$ all-241 $>$ 37-slice
for promoter/fluorescence) --- the evidence base for family-balanced scheduling and
per-family checkpoints (\cref{sec:trainproto}). Per-family, per-task and prior-best
numbers: \cref{tab:leaderboard,tab:priorbest} and the released CSV.
